% Vietnamese translation for OI Public Library
% Source-File: IOI中国国家候选队论文/2005/胡伟栋/胡伟栋.ppt
\documentclass[11pt]{article}
\usepackage{oipl-vn}

\title{Bản trình chiếu: Thuật toán không hoàn hảo}
\author{Hu Weidong}
\date{2005}

\begin{document}
\maketitle

\origtitle{Slide companion for imperfect algorithms}
\sourcepath{IOI China candidate team papers / 2005 / Hu Weidong / Hu Weidong.ppt}

\section{Phạm vi bản dịch}

Bản trình chiếu đi cùng bài viết về thuật toán không hoàn hảo. Các mốc đọc
được gồm \(O(n^{0.5})\), kiểm tra số nguyên tố kiểu \(2^s d\), các số 2047,
1010, 4842, NOIP 2003, IOI 2004 \textit{Polygon}, và Monte Carlo.

\section{Tư tưởng}

Thuật toán không hoàn hảo chấp nhận một mức rủi ro hoặc một giới hạn về tính
tổng quát để đổi lấy tốc độ, bộ nhớ hoặc độ đơn giản. Trong thi lập trình,
cách này hữu ích khi bài có dữ liệu đặc biệt, khi chỉ cần xác suất sai rất
nhỏ, hoặc khi lời giải chính xác quá phức tạp so với thời gian cài đặt.

\section{Ví dụ}

Các slide nhắc tới kiểm tra nguyên tố xác suất: thay vì phân tích thừa số đầy
đủ, ta dùng một số cơ sở nhỏ như \(2,3,5,7\) để loại phần lớn hợp số. Với
\textit{Polygon}, các kỹ thuật hình học hoặc ngẫu nhiên hóa có thể giúp xử lý
các trường hợp khó trong thực nghiệm.

\section{Kết luận}

Điều quan trọng là hiểu rõ rủi ro. Một thuật toán Monte Carlo hoặc heuristic
chỉ nên dùng khi xác suất sai được kiểm soát hoặc khi bối cảnh bài thi cho
phép cách tiếp cận ấy.

\end{document}
